\section{Related Work}
\label{sec:related-work}

\todoconfirm{This section is drafted against general background rather
than the two reference papers Xuanhao sent for how experimental results
should be presented — those weren't available in the session that drafted
this paper. Once shared, this section's structure/tone and citation set
should be revised to match them; the citations below are limited to works
this project directly builds on, to avoid citing anything not actually
verified.}

\paragraph{Orbital edge computing and LEO satellite networking.} Moving
computation to the network edge is a well-established idea in terrestrial
systems; OEC applies the same principle where the "edge" is the satellite
itself, motivated by a downlink budget that is orders of magnitude smaller
than the volume of imagery a modern sensor can capture in a single pass.
LEO constellations add a networking wrinkle largely absent from
terrestrial edge computing: topology and link capacity are not merely
variable but \emph{predictable} — satellite and ground-station positions,
and therefore contact windows and candidate routes, are known in advance
from orbital mechanics. This paper's simulation environment is built on
Hypatia~\cite{kassing2020hypatia}, an open-source LEO network simulator
that models constellation geometry, inter-satellite links, and
ground-station contact windows explicitly, which is what makes a
rolling-horizon scheduler that \emph{predicts} short-term future
connectivity (rather than only reacting to it) both possible and testable
at scale.

\paragraph{Learned image compression for downlink-constrained systems.}
On-board compression here uses a residual-quantized variational
autoencoder (RQ-VAE)~\cite{lee2022autoregressive}: an encoder maps an
image to a small set of codebook indices, and successive
\emph{residual} quantization stages progressively refine the
reconstruction, each stage correcting the error left by the stage before
it. The number of stages run before transmission — this paper's
\emph{depth} — is exactly the fidelity/payload knob a scheduler has to
reason about. We evaluate this compression model on two remote-sensing
benchmarks: EuroSAT~\cite{helber2019eurosat}, a 10-class land-cover
classification dataset used for an initial spatial-size/depth sweep, and
FLAIR-1~\cite{garioud2022flair}, a larger aerial-imagery semantic-segmentation
dataset (15 land-cover classes) used for the paper's central downstream-task
evaluation. Unlike most learned-compression evaluations, which report
pixel-fidelity metrics (PSNR, SSIM, LPIPS, FID) as an end in themselves,
this paper treats those metrics as necessary but insufficient: our results
(\S\ref{sec:downstream}) show they can be nearly flat exactly where a
downstream task's performance is not.

\paragraph{MPC-based scheduling and network optimization.} Model
Predictive Control — solve an optimization over a finite look-ahead
horizon, act on only the first step, then re-plan as new information
arrives — is a natural fit for a system whose future state (contact
windows, link capacities) is partially predictable but subject to real
uncertainty (bursty task arrival, imperfect topology forecasts). We
formulate the scheduling problem as a mixed-integer linear program so it
remains tractable at the problem sizes this paper evaluates, and construct
a piecewise-linear coverage reward specifically so that a concave,
diminishing-returns objective can be optimized without introducing
additional binary or SOS2 variables (\S\ref{sec:coverage}) — a standard
technique for keeping a naturally nonlinear objective inside a linear
program's reach. We evaluate this rolling-horizon approach against both
simpler greedy baselines and a learned (reinforcement-learning) policy
(\S\ref{sec:rl}), and against hierarchical and peer couplings that split
routing and depth decisions into separate, communicating optimizers
(\S\ref{sec:formulation}).
